% ============================================================
% Question Bank: ch09-03 Experimental Probability (Modified — FL/TX simulation emphasis)
% Topic Title: Experimental Probability (Simulation Emphasis)
% CCSS: 7.SP.C.7 (FL/TX modification)
% Questions: 30 (18 MC + 7 SA + 5 Graphical)
% ============================================================

% --- Multiple Choice Questions (18) ---

\begin{practiceQuestion}{m-9-3-q01}{mc}
\begin{questionText}
A coin is flipped $100$ times. Heads appears $48$ times. What is the experimental probability of heads?
\end{questionText}
\choiceA{$0.50$}
\choiceB{$0.48$}
\choiceC{$0.52$}
\choiceD{$0.02$}
\correctAnswer{B}
\explanation{$P = \frac{48}{100} = 0.48$.}
\end{practiceQuestion}

\begin{practiceQuestion}{m-9-3-q02}{mc}
\begin{questionText}
Which of the following BEST describes a simulation?
\end{questionText}
\choiceA{An exact calculation of probability}
\choiceB{A model that uses random tools to imitate a real-world situation}
\choiceC{A graph showing all outcomes}
\choiceD{A formula for finding theoretical probability}
\correctAnswer{B}
\explanation{A simulation uses coins, dice, spinners, or random-number generators to model real events and estimate probabilities.}
\end{practiceQuestion}

\begin{practiceQuestion}{m-9-3-q03}{mc}
\begin{questionText}
A spinner has $4$ equal sections. After $80$ spins, red appears $24$ times. The theoretical probability of red is $0.25$. What is the experimental probability?
\end{questionText}
\choiceA{$0.25$}
\choiceB{$0.30$}
\choiceC{$0.20$}
\choiceD{$0.24$}
\correctAnswer{B}
\explanation{$P = \frac{24}{80} = 0.30$. This is slightly above the theoretical $0.25$, but the difference could be due to chance.}
\end{practiceQuestion}

\begin{practiceQuestion}{m-9-3-q04}{mc}
\begin{questionText}
To design a simulation, what is the FIRST step?
\end{questionText}
\choiceA{Record results in a frequency table}
\choiceB{Run many trials}
\choiceC{Identify the event and its possible outcomes}
\choiceD{Compare experimental and theoretical probabilities}
\correctAnswer{C}
\explanation{The first step is to identify what you are simulating and the possible outcomes before choosing a model.}
\end{practiceQuestion}

\begin{practiceQuestion}{m-9-3-q05}{mc}
\begin{questionText}
You want to simulate a $30\%$ chance of an event using random digits $0$--$9$. Which digits should represent the event?
\end{questionText}
\choiceA{$0, 1, 2$}
\choiceB{$0, 1, 2, 3$}
\choiceC{$0, 1$}
\choiceD{$0, 1, 2, 3, 4$}
\correctAnswer{A}
\explanation{$30\% = \frac{3}{10}$. Using $3$ digits ($0, 1, 2$) out of $10$ gives a $30\%$ probability.}
\end{practiceQuestion}

\begin{practiceQuestion}{m-9-3-q06}{mc}
\begin{questionText}
A die is rolled $60$ times. A $6$ appears $15$ times. What is the experimental probability, and how does it compare to theoretical?
\end{questionText}
\choiceA{Experimental $= 0.25$, higher than theoretical $\frac{1}{6} \approx 0.167$}
\choiceB{Experimental $= 0.167$, equal to theoretical}
\choiceC{Experimental $= 0.15$, lower than theoretical}
\choiceD{Experimental $= 0.25$, lower than theoretical}
\correctAnswer{A}
\explanation{$P = \frac{15}{60} = 0.25$. Theoretical $= \frac{1}{6} \approx 0.167$. The experimental value is noticeably higher.}
\end{practiceQuestion}

\begin{practiceQuestion}{m-9-3-q07}{mc}
\begin{questionText}
A random-number generator produces digits $0$--$9$. To simulate a $40\%$ chance of rain, you let digits $0$--$3$ represent rain. In $10$ trials with digits $3, 7, 1, 0, 9, 4, 2, 8, 5, 6$, how many rain events occur?
\end{questionText}
\choiceA{$3$}
\choiceB{$4$}
\choiceC{$5$}
\choiceD{$6$}
\correctAnswer{B}
\explanation{Rain digits ($0$--$3$): $3, 1, 0, 2$ — that is $4$ rain events out of $10$ trials.}
\end{practiceQuestion}

\begin{practiceQuestion}{m-9-3-q08}{mc}
\begin{questionText}
Which tool would best simulate an event with a $\frac{1}{6}$ probability?
\end{questionText}
\choiceA{A fair coin}
\choiceB{A standard number cube}
\choiceC{A spinner with $4$ equal sections}
\choiceD{A bag with $3$ red and $3$ blue marbles}
\correctAnswer{B}
\explanation{A standard die has $6$ equally likely outcomes, each with probability $\frac{1}{6}$.}
\end{practiceQuestion}

\begin{practiceQuestion}{m-9-3-q09}{mc}
\begin{questionText}
About $30\%$ of blood donors have type~A blood. Using digits $0$--$9$ with $0, 1, 2$ representing type~A, a student runs $5$ trials: $7, 4, 1, 0, 9$. How many trials resulted in type~A?
\end{questionText}
\choiceA{$1$}
\choiceB{$2$}
\choiceC{$3$}
\choiceD{$4$}
\correctAnswer{B}
\explanation{Digits $0$--$2$: only $1$ and $0$ match. That is $2$ trials with type~A.}
\end{practiceQuestion}

\begin{practiceQuestion}{m-9-3-q10}{mc}
\begin{questionText}
A basketball player makes $70\%$ of her free throws. Which simulation model would work?
\end{questionText}
\choiceA{Flip a coin — heads = make}
\choiceB{Roll a die — $1$ or $2$ = make}
\choiceC{Use digits $0$--$9$ — digits $0$--$6$ = make}
\choiceD{Use digits $0$--$9$ — digits $0$--$2$ = make}
\correctAnswer{C}
\explanation{$70\% = \frac{7}{10}$. Digits $0, 1, 2, 3, 4, 5, 6$ ($7$ digits) represent making the shot.}
\end{practiceQuestion}

\begin{practiceQuestion}{m-9-3-q11}{mc}
\begin{questionText}
As the number of trials increases in a simulation, the experimental probability:
\end{questionText}
\choiceA{Moves farther from the theoretical probability}
\choiceB{Stays the same after $10$ trials}
\choiceC{Tends to get closer to the theoretical probability}
\choiceD{Always equals the theoretical probability}
\correctAnswer{C}
\explanation{The Law of Large Numbers tells us that more trials lead to experimental probabilities closer to theoretical values.}
\end{practiceQuestion}

\begin{practiceQuestion}{m-9-3-q12}{mc}
\begin{questionText}
A simulation uses a coin to model whether a family has a boy or girl. Heads = boy, tails = girl. To simulate a family with $3$ children, what do you do?
\end{questionText}
\choiceA{Flip the coin $1$ time}
\choiceB{Flip the coin $3$ times and record the sequence}
\choiceC{Flip the coin until you get heads}
\choiceD{Roll a die $3$ times}
\correctAnswer{B}
\explanation{Each child is equally likely boy or girl (like a coin flip). Flip $3$ times to simulate $3$ children.}
\end{practiceQuestion}

\begin{practiceQuestion}{m-9-3-q13}{mc}
\begin{questionText}
A student runs a simulation $10$ times and a separate simulation $100$ times for the same event. Which simulation gives a more reliable estimate?
\end{questionText}
\choiceA{The $10$-trial simulation, because it is faster}
\choiceB{The $100$-trial simulation, because more data gives a better estimate}
\choiceC{Both are equally reliable}
\choiceD{Neither is reliable}
\correctAnswer{B}
\explanation{More trials provide better accuracy. With $100$ trials, the experimental probability is closer to the true value.}
\end{practiceQuestion}

\begin{practiceQuestion}{m-9-3-q14}{mc}
\begin{questionText}
A frequency table from a simulation shows that "success" occurred $35$ times in $100$ trials. What is the experimental probability of success?
\end{questionText}
\choiceA{$0.035$}
\choiceB{$0.35$}
\choiceC{$3.5$}
\choiceD{$35$}
\correctAnswer{B}
\explanation{$P = \frac{35}{100} = 0.35$.}
\end{practiceQuestion}

\begin{practiceQuestion}{m-9-3-q15}{mc}
\begin{questionText}
Why would you use a simulation instead of a real experiment?
\end{questionText}
\choiceA{Simulations are always $100\%$ accurate}
\choiceB{Real experiments are never useful}
\choiceC{When real experiments are too expensive, slow, or impractical}
\choiceD{Simulations don't require any planning}
\correctAnswer{C}
\explanation{Simulations are used when conducting the real experiment would be too costly, time-consuming, or impractical.}
\end{practiceQuestion}

\begin{practiceQuestion}{m-9-3-q16}{mc}
\begin{questionText}
A student simulates flipping a coin $50$ times. She gets $30$ heads. How does the experimental probability compare to theoretical?
\end{questionText}
\choiceA{Experimental ($0.60$) is higher than theoretical ($0.50$)}
\choiceB{Experimental ($0.50$) equals theoretical ($0.50$)}
\choiceC{Experimental ($0.30$) is lower than theoretical ($0.50$)}
\choiceD{Experimental ($0.60$) equals theoretical ($0.60$)}
\correctAnswer{A}
\explanation{$P = \frac{30}{50} = 0.60$. This is higher than the theoretical $0.50$, but within expected variation for $50$ trials.}
\end{practiceQuestion}

\begin{practiceQuestion}{m-9-3-q17}{mc}
\begin{questionText}
A simulation uses random digits $0$--$9$ to model an event. Digits $0$--$4$ represent success ($50\%$). In $20$ trials, success occurs $8$ times. What is the experimental probability?
\end{questionText}
\choiceA{$0.50$}
\choiceB{$0.40$}
\choiceC{$0.80$}
\choiceD{$0.08$}
\correctAnswer{B}
\explanation{$P = \frac{8}{20} = 0.40$. The theoretical is $0.50$; the difference is expected with only $20$ trials.}
\end{practiceQuestion}

\begin{practiceQuestion}{m-9-3-q18}{mc}
\begin{questionText}
Which of the following is the LAST step in designing a simulation?
\end{questionText}
\choiceA{Choose a model}
\choiceB{Identify the event}
\choiceC{Run many trials}
\choiceD{Calculate the experimental probability and compare to the theoretical value}
\correctAnswer{D}
\explanation{After identifying the event, choosing a model, and running trials, the last step is to calculate the experimental probability and compare it to the theoretical value.}
\end{practiceQuestion}

% --- Short Answer Questions (7) ---

\begin{practiceQuestion}{m-9-3-q19}{sa}
\begin{questionText}
A coin is tossed $200$ times and tails appears $108$ times. What is the experimental probability of tails? Write your answer as a decimal.
\end{questionText}
\correctAnswer{$0.54$}
\explanation{$P = \frac{108}{200} = 0.54$.}
\end{practiceQuestion}

\begin{practiceQuestion}{m-9-3-q20}{sa}
\begin{questionText}
Describe a simulation model to estimate the probability that a family with $4$ children has exactly $2$ boys. Assume each child is equally likely to be a boy or girl.
\end{questionText}
\correctAnswer{Use $4$ coin flips per trial. Heads = boy, tails = girl. Count trials with exactly $2$ heads. Repeat $50$+ trials.}
\explanation{Each child has a $50\%$ chance, modeled by a coin flip. Flip $4$ coins per trial, record how many give exactly $2$ heads, then divide by total trials.}
\end{practiceQuestion}

\begin{practiceQuestion}{m-9-3-q21}{sa}
\begin{questionText}
A random-number generator produces digits $0$--$9$. You want to simulate a $60\%$ probability. Which digits represent the event?
\end{questionText}
\correctAnswer{Digits $0, 1, 2, 3, 4, 5$ (any $6$ of the $10$ digits)}
\explanation{$60\% = \frac{6}{10}$. Choose $6$ digits out of $10$ to represent the event.}
\end{practiceQuestion}

\begin{practiceQuestion}{m-9-3-q22}{sa}
\begin{questionText}
A spinner is spun $150$ times. Red appears $42$ times. What is the experimental probability of landing on red? Write your answer as a decimal.
\end{questionText}
\correctAnswer{$0.28$}
\explanation{$P = \frac{42}{150} = 0.28$.}
\end{practiceQuestion}

\begin{practiceQuestion}{m-9-3-q23}{sa}
\begin{questionText}
A simulation of $80$ trials models an event with theoretical probability $\frac{1}{4}$. How many times would you EXPECT the event to occur?
\end{questionText}
\correctAnswer{$20$}
\explanation{Expected $= 80 \times \frac{1}{4} = 20$.}
\end{practiceQuestion}

\begin{practiceQuestion}{m-9-3-q24}{sa}
\begin{questionText}
A quality inspector tests $500$ items and finds $15$ defective. Based on this simulation, predict how many defective items would be in a shipment of $3{,}000$.
\end{questionText}
\correctAnswer{$90$}
\explanation{$P(\text{defective}) = \frac{15}{500} = 0.03$. In $3{,}000$ items: $0.03 \times 3{,}000 = 90$.}
\end{practiceQuestion}

\begin{practiceQuestion}{m-9-3-q25}{sa}
\begin{questionText}
A simulation shows that in $40$ trials, an event occurs $14$ times. The theoretical probability is $0.30$. Find the experimental probability and determine if it is close to theoretical. Write your answer as a decimal.
\end{questionText}
\correctAnswer{$0.35$; yes, it is close}
\explanation{$P = \frac{14}{40} = 0.35$. The theoretical value is $0.30$. The difference ($0.05$) is small and expected in $40$ trials.}
\end{practiceQuestion}

% --- Graphical Questions (3 GMC + 2 GSA) ---

\begin{practiceQuestion}{m-9-3-q26}{gmc}
\begin{questionText}
The frequency table shows the results of a simulation. A die was rolled $60$ times. Based on the table, what is the experimental probability of rolling a number less than $3$?

\begin{center}
\begin{tabular}{|c|c|}
\hline
\textbf{Number} & \textbf{Frequency} \\
\hline
$1$ & $12$ \\
\hline
$2$ & $8$ \\
\hline
$3$ & $11$ \\
\hline
$4$ & $10$ \\
\hline
$5$ & $9$ \\
\hline
$6$ & $10$ \\
\hline
\end{tabular}
\end{center}
\end{questionText}
\choiceA{$\frac{12}{60}$}
\choiceB{$\frac{20}{60}$}
\choiceC{$\frac{31}{60}$}
\choiceD{$\frac{8}{60}$}
\correctAnswer{B}
\explanation{Numbers less than $3$: $1$ and $2$. Frequency $= 12 + 8 = 20$. $P = \frac{20}{60} = \frac{1}{3}$.}
\end{practiceQuestion}

\begin{practiceQuestion}{m-9-3-q27}{gmc}
\begin{questionText}
The bar graph shows the results of $100$ spins of a spinner that should have equal $25\%$ probability for each color. Which color's experimental probability is farthest from the theoretical $0.25$?

\begin{center}
\begin{tikzpicture}[scale=0.55]
  \draw[thick,->] (0,0) -- (0,38) node[above,font=\small] {Frequency};
  \draw[thick,->] (0,0) -- (9,0);
  \foreach \y in {5,10,...,35} {
    \draw[gray!30] (0,\y) -- (8.5,\y);
    \node[left,font=\small] at (0,\y) {\y};
  }
  \draw[dashed,funRed] (0,25) -- (8.5,25);
  \node[right,font=\small,funRed] at (8.5,25) {Expected};
  \fill[funRed!70] (0.5,0) rectangle (1.5,20);
  \fill[funBlue!70] (2.5,0) rectangle (3.5,28);
  \fill[funGreen!70] (4.5,0) rectangle (5.5,35);
  \fill[funOrange!70] (6.5,0) rectangle (7.5,17);
  \node[below,font=\small] at (1,0) {Red};
  \node[below,font=\small] at (3,0) {Blue};
  \node[below,font=\small] at (5,0) {Grn.};
  \node[below,font=\small] at (7,0) {Org.};
\end{tikzpicture}
\end{center}
\end{questionText}
\choiceA{Red ($20$ times)}
\choiceB{Blue ($28$ times)}
\choiceC{Green ($35$ times)}
\choiceD{Orange ($17$ times)}
\correctAnswer{C}
\explanation{Green appeared $35$ times: $|35 - 25| = 10$. Orange: $|17 - 25| = 8$. Red: $|20 - 25| = 5$. Blue: $|28 - 25| = 3$. Green is farthest from $25$.}
\end{practiceQuestion}

\begin{practiceQuestion}{m-9-3-q28}{gmc}
\begin{questionText}
A student uses random digits $0$--$9$ to simulate a $50\%$ event (digits $0$--$4$ = success). The table shows $10$ trials. How many successes occurred?

\begin{center}
\begin{tabular}{|c|c|c|c|c|c|c|c|c|c|c|}
\hline
\textbf{Trial} & $1$ & $2$ & $3$ & $4$ & $5$ & $6$ & $7$ & $8$ & $9$ & $10$ \\
\hline
\textbf{Digit} & $7$ & $2$ & $5$ & $0$ & $8$ & $4$ & $1$ & $9$ & $3$ & $6$ \\
\hline
\end{tabular}
\end{center}
\end{questionText}
\choiceA{$3$}
\choiceB{$4$}
\choiceC{$5$}
\choiceD{$6$}
\correctAnswer{C}
\explanation{Digits $0$--$4$: trials $2(2), 4(0), 6(4), 7(1), 9(3)$ — that is $5$ successes.}
\end{practiceQuestion}

\begin{practiceQuestion}{m-9-3-q29}{gsa}
\begin{questionText}
The line graph shows the experimental probability of rolling a $1$ on a die after $10$, $20$, $50$, $100$, and $200$ rolls. What is the theoretical probability, and does the graph support the idea that more trials give a better estimate?

\begin{center}
\begin{tikzpicture}[scale=0.7]
  \draw[thick,->] (0,0) -- (6.5,0) node[right,font=\small] {Rolls};
  \draw[thick,->] (0,0) -- (0,5) node[above,font=\small] {$P$(rolling 1)};
  \draw[dashed,gray] (0,2.5) -- (6,2.5);
  \node[left,font=\small] at (0,0) {$0$};
  \node[left,font=\small] at (0,1.25) {$0.08$};
  \node[left,font=\small] at (0,2.5) {$0.167$};
  \node[left,font=\small] at (0,3.75) {$0.25$};
  \node[left,font=\small] at (0,5) {$0.33$};
  \node[below,font=\small] at (1,0) {$10$};
  \node[below,font=\small] at (2,0) {$20$};
  \node[below,font=\small] at (3,0) {$50$};
  \node[below,font=\small] at (4,0) {$100$};
  \node[below,font=\small] at (5,0) {$200$};
  \fill[funBlue] (1,3.75) circle (0.08);
  \fill[funBlue] (2,1.5) circle (0.08);
  \fill[funBlue] (3,3.0) circle (0.08);
  \fill[funBlue] (4,2.7) circle (0.08);
  \fill[funBlue] (5,2.55) circle (0.08);
  \draw[thick,funBlue] (1,3.75) -- (2,1.5) -- (3,3.0) -- (4,2.7) -- (5,2.55);
\end{tikzpicture}
\end{center}
\end{questionText}
\correctAnswer{Theoretical probability is $\frac{1}{6} \approx 0.167$. Yes, the graph shows the experimental probability getting closer to $0.167$ as trials increase.}
\explanation{The dashed line at $0.167$ shows the theoretical value. Early trials ($10, 20$) vary widely, but by $200$ rolls the probability is near $0.167$. This supports the Law of Large Numbers.}
\end{practiceQuestion}

\begin{practiceQuestion}{m-9-3-q30}{gsa}
\begin{questionText}
The frequency table shows two separate simulations of the same event (theoretical probability $= 0.25$). Which simulation gives a more reliable estimate, and why?

\begin{center}
\begin{tabular}{|c|c|c|c|}
\hline
\textbf{Simulation} & \textbf{Trials} & \textbf{Successes} & \textbf{Exp.\ Probability} \\
\hline
A & $20$ & $7$ & $0.35$ \\
\hline
B & $200$ & $52$ & $0.26$ \\
\hline
\end{tabular}
\end{center}
\end{questionText}
\correctAnswer{Simulation B is more reliable because it has more trials ($200$ vs $20$) and its experimental probability ($0.26$) is closer to the theoretical ($0.25$).}
\explanation{More trials reduce the effect of chance variation. Simulation B ($200$ trials, $P = 0.26$) is closer to $0.25$ than Simulation A ($20$ trials, $P = 0.35$).}
\end{practiceQuestion}
